\paperchapter{November 2025}{16 November 2025}{A}{50}{19th GSET}
\keytable{\kcell{1}{A} & \kcell{2}{B} & \kcell{3}{D} & \kcell{4}{D} & \kcell{5}{A} & \kcell{6}{C} & \kcell{7}{C} & \kcell{8}{A} & \kcell{9}{A} & \kcell{10}{D} \\
\kcell{11}{A} & \kcell{12}{A} & \kcell{13}{C} & \kcell{14}{B} & \kcell{15}{B} & \kcell{16}{B} & \kcell{17}{B} & \kcell{18}{B} & \kcell{19}{A} & \kcell{20}{C} \\
\kcell{21}{B} & \kcell{22}{C} & \kcell{23}{A} & \kcell{24}{C} & \kcell{25}{C} & \kcell{26}{D} & \kcell{27}{D} & \kcell{28}{C} & \kcell{29}{B} & \kcell{30}{C} \\
\kcell{31}{D} & \kcell{32}{B} & \kcell{33}{C} & \kcell{34}{C} & \kcell{35}{A} & \kcell{36}{C} & \kcell{37}{A} & \kcell{38}{A} & \kcell{39}{D} & \kcell{40}{C} \\
\kcell{41}{C} & \kcell{42}{A} & \kcell{43}{D} & \kcell{44}{A} & \kcell{45}{C} & \kcell{46}{B} & \kcell{47}{B} & \kcell{48}{B} & \kcell{49}{B} & \kcell{50}{C}}
\begin{question}{1}
{Which of the following is an example of synchronous communication in ICT?\par}
\optstwo{Video conferencing}{E-mails}{Blogs}{Wikis}
\answer{A}{Video conferencing}
\begin{explain}
In \textbf{synchronous} communication all participants are present \emph{at the same time} and interact in real time. Video conferencing (Zoom, Google Meet, Microsoft Teams), live chat, telephone calls and live webinars are synchronous.
\why{E-mails, blogs and wikis are \textbf{asynchronous}: the sender writes now and the receiver reads or edits later, at a time of their own choosing.}
\tip{Synchronous = same time (live). Asynchronous = different time (stored and forwarded).}
\end{explain}
\end{question}

\begin{question}{2}
{Who is known as the Inventor of Artificial Intelligence?\par}
\optstwo{Charles Babbage}{John McCarthy}{Andrew Ng}{Mark Zuckerberg}
\answer{B}{John McCarthy}
\begin{explain}
\textbf{John McCarthy} (1927--2011) coined the term \emph{Artificial Intelligence} in the 1955 proposal for the \textbf{Dartmouth Summer Research Project (1956)}, the event regarded as the birth of AI as a field. He also created the \textbf{LISP} programming language (1958), long the main language of AI research.
\why{Charles Babbage designed the Analytical Engine and is called the ``father of the computer''. Andrew Ng is a modern machine-learning researcher (co-founder of Coursera and Google Brain). Mark Zuckerberg founded Facebook.}
\end{explain}
\end{question}

\begin{question}{3}
{Identify the incorrect pair of Swayam Prabha DTH Channel and its theme/subject.\par}
\optsone{Channel 01 VAGEESH : Language and Literature}{Channel 03 PRABODH : Social and Behavioral Science}{Channel 06 VIDHIK : Law, Legal Studies and Human Rights}{Channel 10 DAKSH : Life Sciences}
\answer{D}{Channel 10 DAKSH : Life Sciences}
\begin{explain}
The 40 SWAYAM Prabha DTH channels are each dedicated to a theme. Channel~01 \emph{Vageesh} (Language \& Literature), Channel~03 \emph{Prabodh} (Social \& Behavioural Sciences) and Channel~06 \emph{Vidhik} (Law, Legal Studies \& Human Rights) are correctly matched. Channel~10 \emph{Daksh} is devoted to \emph{Applied Sciences}; Life Sciences is carried by Channel~09 \emph{Spandan}. Hence pair (D) is incorrect.
\tip{SWAYAM Prabha: 40 free 24$\times$7 educational DTH channels via GSAT-15, launched 9 July 2017; content from CEC--UGC, NPTEL (IITs), IGNOU, NCERT, NIOS and others.}
\end{explain}
\end{question}

\begin{question}{4}
{Which of the following softwares are used in India for plagiarism checking?\newline I. Turnitin\newline II. Drill Bit\newline III. Urkund\newline Choose the correct option:\par}
\optstwo{I and II only}{I and III only}{II and III only}{I, II and III}
\answer{D}{I, II and III}
\begin{explain}
All three are plagiarism-detection tools used by Indian institutions:
\begin{itemize}
\item \textbf{Turnitin} -- widely used international service.
\item \textbf{DrillBit} -- Indian-developed software; provided free to universities by the Ministry of Education/INFLIBNET under the \emph{ShodhShuddhi} initiative (from 2022).
\item \textbf{Urkund} (now \emph{Ouriginal}) -- the tool originally supplied by INFLIBNET under ShodhShuddhi (2019).
\end{itemize}
The UGC (Promotion of Academic Integrity and Prevention of Plagiarism in HEIs) Regulations, 2018 require such checks for theses and dissertations. Hence I, II and III.
\end{explain}
\end{question}

\begin{question}{5}
{Put the following units of storage into the correct order, starting with the largest unit first and going down to the smallest unit.\newline i. Kilobyte\newline ii. Byte\newline iii. Megabyte\newline iv. Terabyte\newline v. Gigabyte\newline vi. Bit\newline Choose the correct option:\par}
\optstwo{iv, v, iii, i, ii, vi}{v, iv, iii, i, ii, vi}{iv, iii, v, i, ii, vi}{iv, v, i, ii, iii, vi}
\answer{A}{iv, v, iii, i, ii, vi}
\begin{explain}
Storage units from the largest to the smallest:
\[\text{Terabyte (iv)} > \text{Gigabyte (v)} > \text{Megabyte (iii)} > \text{Kilobyte (i)} > \text{Byte (ii)} > \text{Bit (vi)}\]
(1~TB = 1024~GB, 1~GB = 1024~MB, 1~MB = 1024~KB, 1~KB = 1024~bytes, 1~byte = 8~bits.) The order iv, v, iii, i, ii, vi is option (A).
\end{explain}
\end{question}

\begin{question}{6}
{The encoding of the communication message which takes into consideration the circumstances and environment refers to the\par}
\optsfour{Decoding}{Channel}{Context}{Response}
\answer{C}{Context}
\begin{explain}
\textbf{Context} is the setting (physical, social, cultural, psychological and temporal) in which communication takes place. A sender who encodes a message keeping the circumstances and environment in mind is taking the \emph{context} into account; the same words may carry different meanings in different contexts.
\why{Decoding is the receiver interpreting the message; the channel is the medium that carries it; the response is the receiver's reaction or feedback.}
\end{explain}
\end{question}

\begin{question}{7}
{Two students who stand apart from each other, maintaining their territorial boundaries, refers to\par}
\optsfour{Haptics}{Kinesics}{Proxemics}{Chronemics}
\answer{C}{Proxemics}
\begin{explain}
\textbf{Proxemics} (term coined by Edward T.~Hall, 1966) is the study of how people use \emph{space and distance} in communication: intimate, personal, social and public zones. Students keeping a distance to maintain their ``territory'' is a proxemic behaviour.
\why{Haptics = touch; Kinesics = body movements, gestures and facial expressions; Chronemics = use of time.}
\end{explain}
\end{question}

\begin{question}{8}
{The teachers give regular feedback to their institutional heads on student performance through \blank{} of communication.\par}
\optstwo{Upward channel}{Horizontal channel}{Downward channel}{Lateral channel}
\answer{A}{Upward channel}
\begin{explain}
Communication that flows from subordinates to superiors (teacher $\rightarrow$ head of institution) is \textbf{upward communication}. Reports, feedback, suggestions and grievances travel upward.
\why{Downward flows from superior to subordinate (circulars, instructions); horizontal/lateral flows between people at the same level (teacher to teacher).}
\end{explain}
\end{question}

\begin{question}{9}
{Perception barrier in communication occurs because people perceive the\par}
\optsone{Same phenomena from different frame of reference}{Different phenomena from different frame of reference}{Different phenomena from same frame of reference}{Same phenomena from same frame of reference}
\answer{A}{Same phenomena from different frame of reference}
\begin{explain}
A \textbf{perceptual barrier} arises because different people look at the \emph{same} event, message or phenomenon from \emph{different frames of reference} (their experiences, beliefs, culture and expectations) and therefore interpret it differently.
\tip{Frame of reference = the mental ``lens'' formed by one's background. Same reality + different lenses = different perceptions.}
\end{explain}
\end{question}

\begin{question}{10}
{Yellow journalism is\par}
\optstwo{Development Journalism}{Print Journalism}{Investigative Journalism}{Tabloid Journalism}
\answer{D}{Tabloid Journalism}
\begin{explain}
\textbf{Yellow journalism} is sensational, exaggerated and poorly researched reporting meant to attract readers; the term comes from the 1890s circulation war in New York between Pulitzer's \emph{New York World} and Hearst's \emph{New York Journal}. Its modern form is \textbf{tabloid journalism}.
\why{Development journalism reports on development issues; investigative journalism involves in-depth, fact-based inquiry; print journalism is simply the medium.}
\end{explain}
\end{question}

\begin{question}{11}
{Find the odd one out from the following.\newline 1. Doctor\newline 2. Hospital\newline 3. Nurse\newline 4. Patient\par}
\optsfour{2}{3}{1}{4}
\answer{A}{2}
\begin{explain}
Doctor, Nurse and Patient are all \emph{persons} found in a hospital; \textbf{Hospital} (item~2) is the \emph{place} itself. Hence item~2 is the odd one out, which is option (A).
\end{explain}
\end{question}

\begin{question}{12}
{The abhava of thing is known through\par}
\optsfour{Anupalabdhi}{Arthapalthi}{Upasthiti}{Upalabdhi}
\answer{A}{Anupalabdhi}
\begin{explain}
In Bhatta Mimamsa and Advaita Vedanta, \textbf{Anupalabdhi} (non-perception or non-cognition) is the independent \emph{pramana} (means of valid knowledge) through which the \textbf{abhava} (non-existence or absence) of a thing is known. For example, ``there is no pot on the floor'' is known by the non-perception of a pot where it would have been perceived.
\why{Arthapatti is postulation or presumption (e.g.\ ``fat Devadatta does not eat by day, so he eats at night''). Upalabdhi means perception, which is the opposite. Upasthiti is not a pramana.}
\tip{Six pramanas (Mimamsa/Vedanta): Pratyaksha, Anumana, Upamana, Shabda, Arthapatti, Anupalabdhi.}
\end{explain}
\end{question}

\begin{question}{13}
{What are the missing letters in the following series?\newline pqr \_ \_ rsprs \_ \_ spq \_\par}
\optsfour{pqppq}{spqpr}{sqpqr}{sqprq}
\answer{C}{sqpqr}
\begin{explain}
Fill the blanks with \textbf{s, q, p, q, r}:
\[\text{pqr\underline{s}}\ \ \text{\underline{q}rsp}\ \ \text{rs\underline{p}\underline{q}}\ \ \text{spq\underline{r}}\]
Each block of four letters is the previous block shifted one place to the left (pqrs $\rightarrow$ qrsp $\rightarrow$ rspq $\rightarrow$ spqr). The missing letters are s, q, p, q, r, i.e.\ ``sqpqr'', option (C).
\end{explain}
\end{question}

\begin{question}{14}
{If `ideas' related to `brain', then `money' related to \blank{}.\par}
\optsfour{Cash}{Bank}{Loan}{Investment}
\answer{B}{Bank}
\begin{explain}
The relation is \emph{thing : place where it is stored or kept}. Ideas are stored in the brain; money is kept (deposited) in a bank.
\why{Cash is a form of money, a loan is borrowed money, and an investment is a use of money. None of these is the place where money is kept.}
\end{explain}
\end{question}

\begin{question}{15}
{What is the correct sequence of the five membered syllogism (Avayavas) in the Nyaya System?\par}
\optsone{Hetu -- Pratijna -- Upanaya -- Udharana -- Nigamana}{Pratijna -- Hetu -- Udharana -- Upanaya -- Nigamana}{Udharana -- Hetu -- Upanaya -- Nigamana -- Pratijna}{Pratijna -- Udharana -- Hetu -- Nigamana -- Upanaya}
\answer{B}{Pratijna -- Hetu -- Udharana -- Upanaya -- Nigamana}
\begin{explain}
The \textbf{Nyaya five-membered syllogism (Panchavayava)} in its correct order:
\begin{enumerate}
\item \textbf{Pratijna} (proposition): The hill has fire.
\item \textbf{Hetu} (reason): Because it has smoke.
\item \textbf{Udaharana} (example): Wherever there is smoke there is fire, as in a kitchen.
\item \textbf{Upanaya} (application): The hill has smoke, which is invariably accompanied by fire.
\item \textbf{Nigamana} (conclusion): Therefore the hill has fire.
\end{enumerate}
This is option (B).
\end{explain}
\end{question}

\begin{question}{16}
{Which of the following best defines descriptive research?\par}
\optsone{Research focused solely on creating new theories}{Investigation focused on presenting status of a phenomena as they exist}{Study to prove scientific law under controlled situation}{Research designed to change existing conditions}
\answer{B}{Investigation focused on presenting status of a phenomena as they exist}
\begin{explain}
\textbf{Descriptive research} describes ``what is'', meaning the current status of a phenomenon as it naturally exists, without manipulating variables. Surveys, case studies, developmental and correlational studies are descriptive.
\why{Creating new theory is the aim of fundamental or grounded-theory research; proving laws under controlled conditions is experimental research; changing existing conditions is the aim of action research.}
\end{explain}
\end{question}

\begin{question}{17}
{An opinionnaire based study on social issues will have\par}
\optstwo{Historical method}{Survey method}{Experimental method}{Quantitative method}
\answer{B}{Survey method}
\begin{explain}
An \textbf{opinionnaire} (attitude scale or opinion questionnaire) collects the opinions of a large number of respondents on social issues. Collecting such data from a sample to describe a population is the \textbf{survey method}, a form of descriptive research.
\why{The historical method studies past records; the experimental method manipulates variables; ``quantitative'' is a type of data or approach, not a specific method of this kind.}
\end{explain}
\end{question}

\begin{question}{18}
{What is an ethically appropriate method when referring to individual participants in research?\par}
\optsone{List all their names with their responses}{Use code language or anonymous identifiers}{Ask participants to publish their own data}{Replace names with name of data collectors}
\answer{B}{Use code language or anonymous identifiers}
\begin{explain}
Research ethics require \textbf{confidentiality and anonymity}: participants should be referred to by codes, pseudonyms or anonymous identifiers (e.g.\ Respondent R-12, Teacher T3) so that responses cannot be traced back to individuals.
\why{Listing names with responses violates confidentiality; asking participants to publish their own data is irrelevant; replacing names with those of the data collectors is misleading.}
\end{explain}
\end{question}

\begin{question}{19}
{Which of the following set has only qualitative measurement?\par}
\optstwo{Values, attitude, anecdotes}{Capital investment, vision, mission}{Values, performance, profit}{Expenditure, investment, profit}
\answer{A}{Values, attitude, anecdotes}
\begin{explain}
Qualitative measurement deals with \emph{non-numerical} attributes. \textbf{Values, attitudes and anecdotes} are all qualitative, described in words, categories or ratings of quality.
\why{Every other set contains quantitative items: capital investment, profit, performance scores and expenditure are measured in numbers.}
\end{explain}
\end{question}

\begin{question}{20}
{What percentage of observations are broadly covered from `0' SD to `-3' SD of normal probability distribution curve?\par}
\optsfour{100\%}{0\%}{50\%}{30\%}
\answer{C}{50\%}
\begin{explain}
The normal curve is symmetric about the mean (0~SD), so exactly half the observations lie on each side. The area from the mean to $-3\,\sigma$ is 49.87\%, which is broadly \textbf{50\%}.
\begin{center}
\begin{tikzpicture}
\begin{axis}[width=8cm,height=3.6cm,axis lines=middle,ytick=\empty,xtick={-3,-2,-1,0,1,2,3},
  xticklabels={$-3\sigma$,$-2\sigma$,$-1\sigma$,$0$,$1\sigma$,$2\sigma$,$3\sigma$},
  ymin=0,ymax=0.45,xmin=-3.5,xmax=3.5,tick label style={font=\scriptsize}]
\addplot[domain=-3:0,samples=60,fill=brand!25,draw=none] {exp(-x^2/2)/sqrt(2*pi)} \closedcycle;
\addplot[domain=-3.5:3.5,samples=100,thick,brand] {exp(-x^2/2)/sqrt(2*pi)};
\node[font=\scriptsize] at (axis cs:-1.1,0.12) {49.87\%};
\end{axis}
\end{tikzpicture}
\end{center}
\tip{Mean to $\pm1\sigma$: 34.13\% each side (68.26\% total); $\pm2\sigma$: 95.44\%; $\pm3\sigma$: 99.73\%.}
\end{explain}
\end{question}

\begin{question}{21}
{`Gujarat Public Universities Bill' was passed in Gujarat Legislative Assembly in the year\par}
\optsfour{2022}{2023}{2020}{2021}
\answer{B}{2023}
\begin{explain}
The \textbf{Gujarat Public Universities Bill, 2023} was passed by the Gujarat Legislative Assembly in \textbf{September 2023} and became the \emph{Gujarat Public Universities Act, 2023}. It brings 11 state public universities (e.g.\ Gujarat University, M.S.~University of Baroda, Saurashtra University) under a common Act, replacing their separate Acts, with a Board of Management and a Chancellor appointed by the State.
\end{explain}
\end{question}

\begin{question}{22}
{The primary focus of the National Education Policy -- 2020 is\par}
\optsone{Creating an expensive education system}{Providing free education for all}{Transforming India into a knowledge based society and economy}{Emphasizing traditional education}
\answer{C}{Transforming India into a knowledge based society and economy}
\begin{explain}
NEP~2020 states its vision as an India-centred education system that contributes directly to \textbf{transforming India into an equitable and vibrant knowledge society}, by providing high-quality education to all and making India a global knowledge superpower.
\why{It does not seek an expensive system; ``free education for all'' is not its primary focus (RTE already covers ages 6--14); it blends traditional knowledge with modern education rather than emphasising only tradition.}
\end{explain}
\end{question}

\begin{question}{23}
{Name the parent/route organization for DTH Channel 7 (Kautilya) of Swayam Prabha managed by CEC-UGC, New Delhi.\par}
\optstwo{Gujarat University, EMRC, Ahmedabad}{Punjabi University, EMRC, Patiala}{Calicut University, EMRC, Calicut}{IIT, Delhi}
\answer{A}{Gujarat University, EMRC, Ahmedabad}
\begin{explain}
SWAYAM Prabha channels are produced by partner institutions. Channels coordinated by CEC--UGC are prepared through its \textbf{Educational Multimedia Research Centres (EMRCs)}. Channel~7 (\emph{Kautilya}) is handled by the \textbf{EMRC, Gujarat University, Ahmedabad}.
\tip{CEC (Consortium for Educational Communication) is an inter-university centre of UGC with a network of about 21 EMRCs/AVRCs across India.}
\end{explain}
\end{question}

\begin{question}{24}
{Which University in India is offering the largest number of courses through distance mode?\par}
\optsone{University of Mumbai}{University of Delhi}{Indira Gandhi National Open University, Delhi}{Baba Saheb Ambedkar Open University, Ahmedabad}
\answer{C}{Indira Gandhi National Open University, Delhi}
\begin{explain}
\textbf{IGNOU}, the Indira Gandhi National Open University (est.~1985, New Delhi), is a central open university and the largest in the world by enrolment. It offers the largest number of programmes (over 200) through Open and Distance Learning. It also serves as the national apex body for ODL.
\why{Mumbai and Delhi universities offer distance courses through their distance-education wings, but far fewer. BAOU is Gujarat's state open university.}
\end{explain}
\end{question}

\begin{question}{25}
{Which of the following is a common entrance exam for admission to Medical College in India?\par}
\optsfour{JEE}{CAT}{NEET}{GATE}
\answer{C}{NEET}
\begin{explain}
\textbf{NEET-UG} (National Eligibility-cum-Entrance Test), conducted by the \textbf{National Testing Agency (NTA)}, is the single common entrance examination for MBBS, BDS and AYUSH courses in India.
\why{JEE is for engineering (IITs/NITs); CAT is for management (IIMs); GATE is for postgraduate engineering admissions and PSU recruitment.}
\end{explain}
\end{question}

\begin{question}{26}
{The data on the sales of cellular phones (in '000 Rs.) during 2017 -- 2022 years are 48, 40, 30, 25, 28, 40, respectively. The two years between which the rate of change of sales of cellular phones has been minimum are\par}
\optsfour{2017 -- 2018}{2018 -- 2019}{2019 -- 2020}{2020 -- 2021}
\answer{D}{2020 -- 2021}
\begin{explain}
Year-to-year changes in sales:
\begin{center}\small
\begin{tabular}{lccc}\toprule
Period & From $\rightarrow$ To & Change & \% change\\\midrule
2017--18 & 48 $\rightarrow$ 40 & $-8$ & $-16.7\%$\\
2018--19 & 40 $\rightarrow$ 30 & $-10$ & $-25.0\%$\\
2019--20 & 30 $\rightarrow$ 25 & $-5$ & $-16.7\%$\\
\textbf{2020--21} & \textbf{25 $\rightarrow$ 28} & $\mathbf{+3}$ & $\mathbf{+12.0\%}$\\
2021--22 & 28 $\rightarrow$ 40 & $+12$ & $+42.9\%$\\\bottomrule
\end{tabular}
\end{center}
The smallest change, in both absolute and percentage terms, is between \textbf{2020 and 2021}.
\end{explain}
\end{question}

\begin{question}{27}
{The number of candidates (in lakhs) appearing in an entrance examination from five different locations along with the ratio of passing and failing is given as follows.\par}
\qfigure{\begin{tabular}{ccc}\toprule
\textbf{Location} & \textbf{Candidates (lakh)} & \textbf{Pass : Fail}\\\midrule
P & 1.25 & 7 : 3\\ Q & 3.14 & 5 : 3\\ R & 1.08 & 4 : 5\\ S & 2.27 & 1 : 3\\ T & 1.85 & 3 : 2\\\bottomrule
\end{tabular}}
What is the ratio of the number of candidates failing the exam from location S to that of those failing the exam from location P?
\optsfour{50 : 227}{42 : 289}{289 : 42}{227 : 50}
\answer{D}{227 : 50}
\begin{explain}
Failed from S $= 2.27 \times \dfrac{3}{1+3} = 2.27\times\dfrac34 = 1.7025$ lakh.\\
Failed from P $= 1.25 \times \dfrac{3}{7+3} = 1.25\times\dfrac{3}{10} = 0.375$ lakh.
\[\frac{1.7025}{0.375} = \frac{2.27\times 3/4}{1.25\times 3/10} = \frac{2.27\times 10}{1.25\times 4} = \frac{22.7}{5} = \frac{227}{50}\]
Ratio S : P $=$ \textbf{227 : 50}.
\end{explain}
\end{question}

\begin{question}{28}
{The chance that doctor A will diagnose a disease X correctly is 60 percent. The chance that a patient will die by his treatment after correct diagnosis is 40 percent, and the chance of death by wrong diagnosis is 70 percent. A patient of doctor A, who had disease X died. What is the probability that his disease was diagnosed correctly?\par}
\optsfour{0.28}{0.39}{0.46}{0.52}
\answer{C}{0.46}
\begin{explain}
This is Bayes' theorem. Let $C$ = correct diagnosis and $D$ = death.
\[P(C)=0.6,\quad P(D\,|\,C)=0.4;\qquad P(\bar C)=0.4,\quad P(D\,|\,\bar C)=0.7\]
\[P(C\,|\,D)=\frac{0.6\times0.4}{0.6\times0.4+0.4\times0.7}=\frac{0.24}{0.24+0.28}=\frac{0.24}{0.52}\approx \mathbf{0.46}\]
\begin{center}
\begin{tikzpicture}[font=\scriptsize,>=Stealth,node distance=4mm]
\node (p) {Patient with X};
\node (c) at (3,0.6) {Correct (0.6)}; \node (w) at (3,-0.6) {Wrong (0.4)};
\node (cd) at (6.2,0.6) {Dies (0.4): $0.24$}; \node (wd) at (6.2,-0.6) {Dies (0.7): $0.28$};
\draw[->] (p.east) -- (c.west); \draw[->] (p.east) -- (w.west);
\draw[->] (c) -- (cd); \draw[->] (w) -- (wd);
\end{tikzpicture}
\end{center}
\end{explain}
\end{question}

\begin{question}{29}
{The probability that a student is not a swimmer is 1/5. Out of 5 students, what is the probability that atleast 4 are swimmers?\par}
\optsfour{0.65}{0.74}{0.79}{0.82}
\answer{B}{0.74}
\begin{explain}
$P(\text{not swimmer})=\tfrac15 \Rightarrow p=P(\text{swimmer})=\tfrac45=0.8,\ q=0.2$, $n=5$ (binomial).
\[P(X\ge4)=\binom54(0.8)^4(0.2)+\binom55(0.8)^5=5(0.4096)(0.2)+0.32768=0.4096+0.32768=0.73728\]
$\approx$ \textbf{0.74}.
\end{explain}
\end{question}

\begin{question}{30}
{Some tickets are numbered from 1, 2, \ldots{}, 100 for a random draw. What is the probability that the number drawn is a multiple of 3 or 5?\par}
\optsfour{0.27}{0.37}{0.47}{0.57}
\answer{C}{0.47}
\begin{explain}
Multiples of 3 in 1--100: $\lfloor100/3\rfloor=33$. Multiples of 5: $20$. Multiples of both (i.e.\ of 15): $6$.
\[n(3\cup5)=33+20-6=47\qquad\Rightarrow\qquad P=\frac{47}{100}=\mathbf{0.47}\]
\end{explain}
\end{question}

\begin{question}{31}
{Which of the following is not a state of memory level?\par}
\optsfour{Retention}{Recall}{Recognition}{Assimilation}
\answer{D}{Assimilation}
\begin{explain}
Memory is studied through the processes of \textbf{encoding/learning, retention (storage), recall and recognition}, and these are the recognised states or levels at which memory is measured. \textbf{Assimilation} is a Piagetian concept of \emph{cognitive development} (fitting new experience into existing schemas). It is not a state of memory.
\end{explain}
\end{question}

\begin{question}{32}
{For teachers, it is important to consider the individual differences amongst learners. This helps to\par}
\optsone{Classify the students in to subgroups for different activities}{Guide learning and develop students' cognitive processes}{Ensure uniform teaching for all students}{Reduce the need for assessment}
\answer{B}{Guide learning and develop students' cognitive processes}
\begin{explain}
Recognising individual differences (in ability, pace, learning style, interest and background) lets a teacher \textbf{guide learning} appropriately and \textbf{develop each student's cognitive processes} through differentiated instruction.
\why{Classifying students into subgroups is at best a tool, not the purpose; uniform teaching ignores differences; individual differences do not reduce the need for assessment.}
\end{explain}
\end{question}

\begin{question}{33}
{Which of the following is most significant for student learning?\par}
\optstwo{Student background}{Prior academic achievement}{Active learning processes}{Availability of textbooks}
\answer{C}{Active learning processes}
\begin{explain}
Modern learning theory (constructivism) holds that learning happens when learners \textbf{actively} construct knowledge by doing, discussing, questioning and reflecting. Active learning processes therefore matter most, while background, prior achievement and textbooks are supporting factors.
\end{explain}
\end{question}

\begin{question}{34}
{What type of content is broadcast on the Swayam Prabha Channel?\par}
\optsone{Weather and entertainment programmes}{Political debate programmes}{Educational programmes covering diverse fields for the 18+ age group learners}{Telecast of job opportunities to apply}
\answer{C}{Educational programmes covering diverse fields for the 18+ age group learners}
\begin{explain}
SWAYAM Prabha telecasts high-quality \textbf{educational programmes in diverse disciplines} (arts, science, commerce, engineering, social sciences, law, medicine, agriculture and more) for school, higher-education and lifelong learners. For the higher-education channels the target group is learners aged roughly 18 and above.
\end{explain}
\end{question}

\begin{question}{35}
{Under CBCS system, the requirement for awarding degree is determined by\par}
\optsone{Number of credits earned by the students}{Number of years spent in the programme}{Marks obtained in semester examination}{Attendance and Internal Assessment (IA) marks only}
\answer{A}{Number of credits earned by the students}
\begin{explain}
Under the \textbf{Choice Based Credit System (CBCS)}, introduced by UGC in 2015, each course carries credits, and a degree is awarded when the student earns the \textbf{prescribed minimum number of credits} across core, elective and ability- or skill-enhancement courses. Grades are expressed as SGPA/CGPA.
\why{Years of study, raw marks or attendance alone do not decide the award of the degree under CBCS.}
\end{explain}
\end{question}

\begin{question}{36}
{`Fly ash' produced in thermal power plant is useful in\newline i. Road construction\newline ii. Wasteland development\newline iii. Agriculture as micro-nutrient\newline iv. Brick making\newline Choose the correct option:\par}
\optstwo{i, ii and iii only}{ii, iii and iv only}{i, ii, iii and iv}{i, ii and iv only}
\answer{C}{i, ii, iii and iv}
\begin{explain}
Fly ash, the fine residue of coal combustion, is used for:
\begin{itemize}
\item \textbf{road and embankment construction} (the MoRTH/IRC encourage its use);
\item \textbf{wasteland reclamation and development};
\item \textbf{agriculture}, as a soil conditioner supplying micro-nutrients (B, Mo, Zn, S) and improving water-holding capacity;
\item \textbf{brick making}: fly-ash bricks are mandated near thermal power plants by the MoEF notification of 1999.
\end{itemize}
All four uses are correct, so the answer is (C).
\end{explain}
\end{question}

\begin{question}{37}
{Arrange the following themes of Sustainable Development Goals according to their goal number in decreasing order.\newline I. Life below water\newline II. Clean water and sanitation\newline III. Life on land\newline IV. Climate action\newline V. Affordable and clean energy\newline Choose the correct option:\par}
\optstwo{III, I, IV, V, II}{III, I, V, IV, II}{III, I, II, IV, V}{III, I, IV, II, V}
\answer{A}{III, I, IV, V, II}
\begin{explain}
SDG numbers: Life on land = \textbf{15}, Life below water = \textbf{14}, Climate action = \textbf{13}, Affordable and clean energy = \textbf{7}, Clean water and sanitation = \textbf{6}.
Decreasing order: 15 (III), 14 (I), 13 (IV), 7 (V), 6 (II), giving \textbf{III, I, IV, V, II}.
\tip{SDGs 6--7--13--14--15 = Water, Energy, Climate, Oceans, Land.}
\end{explain}
\end{question}

\begin{question}{38}
{Given below are two statements.\newline Statement -- I : Cyclones developed in the region lying between the Tropic of Capricorn and Cancer are called tropical cyclones.\newline Statement -- II : Hurricane Helene was a large, catastrophic and deadly tropical cyclone that caused widespread destruction and fatalities across the South-Eastern United States in September 2024.\newline In light of the above statements, choose the correct answer from the options given below:\par}
\optsone{Both Statement -- I and Statement -- II are true}{Both Statement -- I and Statement -- II are false}{Statement -- I is true, but Statement -- II is false}{Statement -- I is false, but Statement -- II is true}
\answer{A}{Both Statement -- I and Statement -- II are true}
\begin{explain}
\textbf{Statement I} is true: cyclones that originate over warm tropical oceans, in the belt between the Tropic of Cancer and the Tropic of Capricorn, are called tropical cyclones (hurricanes, typhoons and so on).
\textbf{Statement II} is true: Hurricane Helene (late September 2024) made landfall in Florida's Big Bend as a Category~4 storm and caused catastrophic flooding and more than 200 deaths across the south-eastern US, especially North Carolina.
Both statements are true, so the answer is (A).
\end{explain}
\end{question}

\begin{question}{39}
{How many States in India are identified as highly vulnerable to climate change as per the National Climate Vulnerability Assessment Report by the Department of Science and Technology, Govt. of India (2024)?\par}
\optsfour{5}{6}{7}{8}
\answer{D}{8}
\begin{explain}
The Department of Science \& Technology's national report \emph{Climate Vulnerability Assessment for Adaptation Planning in India using a Common Framework} identified \textbf{eight} eastern and north-eastern states as highly vulnerable: Jharkhand, Mizoram, Odisha, Chhattisgarh, Assam, Bihar, Arunachal Pradesh and West Bengal.
\end{explain}
\end{question}

\begin{question}{40}
{Which of the following oxides of Nitrogen is a constituent of Air Quality Index?\par}
\optstwo{Nitric Oxide (NO)}{Nitrous Oxide (N$_2$O)}{Nitrogen Dioxide (NO$_2$)}{Dinitrogen Tetroxide (N$_2$O$_4$)}
\answer{C}{Nitrogen Dioxide (NO$_2$)}
\begin{explain}
India's \textbf{National Air Quality Index} (launched 2014) uses eight pollutants: PM$_{10}$, PM$_{2.5}$, \textbf{NO$_2$}, SO$_2$, CO, O$_3$, NH$_3$ and Pb. Of the nitrogen oxides, only \textbf{nitrogen dioxide} is an AQI constituent.
\tip{AQI categories: Good (0--50), Satisfactory (51--100), Moderately polluted (101--200), Poor (201--300), Very poor (301--400), Severe (401--500).}
\end{explain}
\end{question}

\begin{question}{41}
{If the first and second letter in the word `COMMUNICATIONS' were interchanged, also the third and the fourth letters, the fifth and the sixth letters and so on, which letter would be the tenth letter counting from your right?\par}
\optsfour{A}{T}{N}{U}
\answer{C}{N}
\begin{explain}
COMMUNICATIONS has 14 letters. Swap each pair:
\begin{center}\ttfamily
\begin{tabular}{*{14}{c}}
C&O&M&M&U&N&I&C&A&T&I&O&N&S\\
O&C&M&M&N&U&C&I&T&A&O&I&S&N\\
\end{tabular}
\end{center}
10th from the right = $(14-10+1)$ = 5th from the left = \textbf{N}.
\end{explain}
\end{question}

\begin{question}{42}
{The angle A of a triangle ABC is equal to the sum of the two other angles. Also the ratio of the angle B to angle C is 4 : 5. The three angles are\par}
\optsfour{90\textdegree{}, 40\textdegree{}, 50\textdegree{}}{90\textdegree{}, 55\textdegree{}, 35\textdegree{}}{90\textdegree{}, 60\textdegree{}, 30\textdegree{}}{90\textdegree{}, 70\textdegree{}, 20\textdegree{}}
\answer{A}{90\textdegree{}, 40\textdegree{}, 50\textdegree{}}
\begin{explain}
$A=B+C$ and $A+B+C=180^\circ \Rightarrow 2A=180^\circ \Rightarrow A=90^\circ$, so $B+C=90^\circ$.
With $B:C=4:5$: $B=\tfrac49\times90^\circ=40^\circ$, $C=\tfrac59\times90^\circ=50^\circ$. The angles are \textbf{90\textdegree, 40\textdegree, 50\textdegree}.
\end{explain}
\end{question}

\begin{question}{43}
{Find the missing number (?) in the following series.\newline 4, 8, 12, 24, 36, 72, (?)\par}
\optsfour{92}{96}{98}{108}
\answer{D}{108}
\begin{explain}
The series alternately multiplies by 2 and by 1.5:
\[4\xrightarrow{\times2}8\xrightarrow{\times1.5}12\xrightarrow{\times2}24\xrightarrow{\times1.5}36\xrightarrow{\times2}72\xrightarrow{\times1.5}\mathbf{108}\]
(Equivalently, the differences run $+4,+4,+12,+12,+36,+36$.)
\end{explain}
\end{question}

\begin{question}{44}
{The monthly income of a person was Rs. 13,500 and his monthly expenditure was Rs. 9,000. Next year, his income increased by 14\% and his expenditure by 7\%. Find the percentage increase in his savings.\par}
\optsfour{28\%}{30\%}{32\%}{34\%}
\answer{A}{28\%}
\begin{explain}
Original saving $=13{,}500-9{,}000=4{,}500$.\\
New income $=13{,}500\times1.14=15{,}390$; new expenditure $=9{,}000\times1.07=9{,}630$.\\
New saving $=15{,}390-9{,}630=5{,}760$.
\[\text{Increase}=\frac{5{,}760-4{,}500}{4{,}500}\times100=\frac{1{,}260}{4{,}500}\times100=\mathbf{28\%}\]
\end{explain}
\end{question}

\begin{question}{45}
{Find the compound interest on Rs. 1,000 at the rate of 10\% per annum for 18 months when interest is compounded half-yearly.\par}
\optsfour{Rs. 145.25}{Rs. 152.75}{Rs. 157.63}{Rs. 161.49}
\answer{C}{Rs. 157.63}
\begin{explain}
Half-yearly compounding: rate $=10/2=5\%$ per half-year, periods $=18$ months $\div 6 = 3$.
\[A=1000(1.05)^3=1000\times1.157625=1157.63\qquad CI=1157.63-1000=\textbf{Rs.\ 157.63}\]
\end{explain}
\end{question}

\begin{context}
\textbf{Directions for Q.~46 to 50:} Read the following paragraph and answer the questions that follow.\par\smallskip
Effective learning requires a comprehensive approach that involves appropriate curriculum, engaging pedagogy, continuous formative assessment, and adequate student support. The curriculum must be interesting and relevant, and updated regularly to align with the latest knowledge requirements and to meet specified learning outcomes. High-quality pedagogy is then necessary to successfully impart the curricular material to students; pedagogical practices determine the learning experiences that are provided to students, thus directly influencing learning outcomes. The assessment methods must be scientific, designed to continuously improve learning and test the application of knowledge.

To promote creativity, institutions and faculty will have the autonomy to innovate on matters of curriculum, pedagogy, and assessment within a broad framework of higher education qualifications that ensures consistency across institutions and programmes and across the ODL, online, and traditional `in-class' modes. All assessment systems shall also be decided by the HEI, including those that lead to final certification. The Choice Based Credit System (CBCS) will be revised for instilling innovation and flexibility. HEIs shall move to a criterion-based grading system that assesses student's achievement based on the learning goals for each programme, making the system fairer and outcomes more comparable.
\end{context}
\begin{question}{46}
{What are the essential components of effective learning?\par}
\optsone{Curriculum, textbooks, examination and evaluation}{Curriculum, pedagogy, assessment and student support}{Teachers, pedagogy, technology and curricular material}{Student support, knowledge application, assessment and syllabus}
\answer{B}{Curriculum, pedagogy, assessment and student support}
\begin{explain}
The opening sentence lists the four components: ``appropriate \textbf{curriculum}, engaging \textbf{pedagogy}, continuous formative \textbf{assessment}, and adequate \textbf{student support}.''
(This passage is taken from NEP~2020, paragraph~12.1.)
\end{explain}
\end{question}

\begin{question}{47}
{How do pedagogical practices influence students?\par}
\optsone{Enhance grades}{Influence learning experiences and outcomes}{Reduce homework}{Improve classroom interactions}
\answer{B}{Influence learning experiences and outcomes}
\begin{explain}
The passage says pedagogical practices ``determine the \textbf{learning experiences} that are provided to students, thus directly influencing \textbf{learning outcomes}.''
\why{Grades, homework and classroom interaction are not what the passage names.}
\end{explain}
\end{question}

\begin{question}{48}
{What type of autonomy will institutions and faculty have to promote creativity?\par}
\optsone{Autonomy in determining student fees}{Autonomy in curriculum, pedagogy and assessment}{Autonomy in teacher recruitment}{Autonomy in time-table planning}
\answer{B}{Autonomy in curriculum, pedagogy and assessment}
\begin{explain}
``To promote creativity, institutions and faculty will have the autonomy to innovate on matters of \textbf{curriculum, pedagogy, and assessment}.''
\end{explain}
\end{question}

\begin{question}{49}
{What will the revised Choice Based Credit System (CBCS) aim to achieve?\par}
\optstwo{Simplify exams}{Encourage innovation}{Shorten the academic year}{Reduce the number of courses}
\answer{B}{Encourage innovation}
\begin{explain}
``The Choice Based Credit System (CBCS) will be revised for instilling \textbf{innovation} and flexibility.''
\end{explain}
\end{question}

\begin{question}{50}
{What grading system will HEIs adopt for final certification?\par}
\optstwo{Percentage based}{Rank based}{Criterion based}{Time based}
\answer{C}{Criterion based}
\begin{explain}
``HEIs shall move to a \textbf{criterion-based grading system} that assesses student's achievement based on the learning goals for each programme.'' Criterion-referenced grading measures attainment against fixed learning goals rather than against other students (norm-referenced grading).
\end{explain}
\end{question}
